\documentclass{article}
\usepackage{amsmath,amssymb,amsthm}
\usepackage{algorithm}
\usepackage[noend]{algpseudocode}
\usepackage{hyperref}
\usepackage{geometry}
\geometry{margin=1in}
\newtheorem{theorem}{Theorem}
\newtheorem{lemma}{Lemma}
\newtheorem{assumption}{Assumption}
\newtheorem{corollary}{Corollary}
\newcommand{\prox}{\operatorname{prox}}
\newcommand{\grad}{\nabla}
\newcommand{\R}{\mathbb{R}}

\begin{document}

\title{Proximal Gradient Method \\ Detailed Algorithmic Description and Convergence Proof}
\author{}
\date{}
\maketitle

\tableofcontents
\newpage

%%%%%%%%%%%%%%%%%%%%%%%%%%%%%%%%%%%%%%%%%%%%%%%%%%%%%%%%%%%%%%%%%%%%%%%%%%%%
\section*{Algorithm Name}
\addcontentsline{toc}{section}{Algorithm Name}
\begin{center}
\textbf{Proximal Gradient Method (Forward–Backward Splitting)}
\end{center}
%%%%%%%%%%%%%%%%%%%%%%%%%%%%%%%%%%%%%%%%%%%%%%%%%%%%%%%%%%%%%%%%%%%%%%%%%%%%

\section*{Mathematical Formulation}
\addcontentsline{toc}{section}{Mathematical Formulation}
We consider the composite convex optimization problem  

\[
\min_{x\in \R^{d}} \; F(x) \quad\text{with}\quad 
F(x):=f(x)+g(x),
\]

where  

\begin{itemize}
\item $f:\R^{d}\rightarrow\R$ is convex, differentiable and has an $L$‑Lipschitz continuous gradient, i.e.  
\[
\|\grad f(x)-\grad f(y)\| \le L\|x-y\|\quad\forall x,y\in\R^{d}.
\]
\item $g:\R^{d}\rightarrow\R\cup\{+\infty\}$ is convex, possibly non‑smooth, but its proximal operator can be evaluated in closed form.
\end{itemize}

For a stepsize $\eta>0$, the **proximal operator** of $\eta g$ is  

\[
\prox_{\eta g}(u):=\arg\min_{x\in\R^{d}}
\Big\{ g(x)+\frac{1}{2\eta}\|x-u\|^{2}\Big\}.
\]

The algorithm generates a sequence $\{x_k\}_{k\ge0}$ by  

\[
\boxed{\,x_{k+1}= \prox_{\eta g}\bigl(x_k-\eta \,\grad f(x_k)\bigr)\,},\qquad k=0,1,2,\dots
\tag{1}
\]

The stepsize is usually chosen as  

\[
0<\eta \le \frac{1}{L}.
\tag{2}
\]

The key property we shall use is the **non‑expansiveness** of the proximal operator:

\[
\|\prox_{\eta g}(u)-\prox_{\eta g}(v)\|\le \|u-v\|\quad\forall u,v\in\R^{d}.
\tag{3}
\]

%%%%%%%%%%%%%%%%%%%%%%%%%%%%%%%%%%%%%%%%%%%%%%%%%%%%%%%%%%%%%%%%%%%%%%%%%%%%
\section*{Pseudocode}
\addcontentsline{toc}{section}{Pseudocode}
\begin{algorithm}[H]
\caption{Proximal Gradient Method}
\begin{algorithmic}[1]
\Require Initial point $x_0\in\R^{d}$, stepsize $\eta\in(0,1/L]$, maximum iterations $K$ (or tolerance $\epsilon$)
\For{$k=0$ \textbf{to} $K-1$}
    \State $u_k \gets x_k - \eta\,\grad f(x_k)$ \hfill \Comment{gradient step}
    \State $x_{k+1} \gets \prox_{\eta g}(u_k)$ \hfill \Comment{proximal step}
    \If{$\|x_{k+1}-x_k\|\le\epsilon$}
        \State \textbf{break}
    \EndIf
\EndFor
\State \Return $x_{K}$ (or last iterate)
\end{algorithmic}
\end{algorithm}
%%%%%%%%%%%%%%%%%%%%%%%%%%%%%%%%%%%%%%%%%%%%%%%%%%%%%%%%%%%%%%%%%%%%%%%%%%%%

\section*{Dry Run Example}
\addcontentsline{toc}{section}{Dry Run Example}
We minimize the scalar composite  

\[
F(w)=f(w)+g(w),\qquad 
f(w)=(w-3)^{2},\qquad g(w)=0,
\]

so the proximal operator reduces to the identity.  
Take $w_0=0$, stepsize $\eta=0.1$, and run three iterations.

\begin{center}
\begin{tabular}{c|c|c|c}
Iteration $k$ & $ \grad f(w_k)$ & $u_k = w_k-\eta\grad f(w_k)$ & $w_{k+1}= \prox_{\eta g}(u_k)=u_k$ \\ \hline
0 & $2(w_0-3)=-6$ & $0-0.1(-6)=0.6$ & $0.6$ \\
1 & $2(0.6-3)=-4.8$ & $0.6-0.1(-4.8)=1.08$ & $1.08$ \\
2 & $2(1.08-3)=-3.84$ & $1.08-0.1(-3.84)=1.464$ & $1.464$
\end{tabular}
\end{center}

The objective values are  

\[
F(w_0)=9,\;F(w_1)= (0.6-3)^2=5.76,\;F(w_2)=(1.08-3)^2=3.6864,\;F(w_3)=(1.464-3)^2=2.365\, .
\]

The iterates clearly move toward the minimizer $w^\star=3$.

%%%%%%%%%%%%%%%%%%%%%%%%%%%%%%%%%%%%%%%%%%%%%%%%%%%%%%%%%%%%%%%%%%%%%%%%%%%%
\section*{Assumptions}
\addcontentsline{toc}{section}{Assumptions}
\begin{assumption}[Convexity and Smoothness of $f$]\label{as:convex_smooth}
$f$ is convex and differentiable on $\R^{d}$ and its gradient is $L$‑Lipschitz, i.e.  

\[
\|\grad f(x)-\grad f(y)\|\le L\|x-y\|,\qquad\forall x,y\in\R^{d}.
\tag{A1}
\]
\end{assumption}

\begin{assumption}[Convexity of $g$]\label{as:convex_g}
$g$ is proper, lower‑semicontinuous and convex (possibly nonsmooth).  
Its proximal operator $\prox_{\eta g}$ is well defined for every $\eta>0$.
\tag{A2}
\end{assumption}

\begin{assumption}[Stepsize]\label{as:stepsize}
The stepsize satisfies  

\[
0<\eta\le\frac{1}{L}.
\tag{A3}
\]
\end{assumption}

\begin{assumption}[Existence of a minimizer]\label{as:optimum}
The set of optimal solutions  

\[
\mathcal X^\star:=\arg\min_{x\in\R^{d}}F(x)
\]

is non‑empty and we denote $F^\star:=\inf_xF(x)$.
\tag{A4}
\end{assumption}

%%%%%%%%%%%%%%%%%%%%%%%%%%%%%%%%%%%%%%%%%%%%%%%%%%%%%%%%%%%%%%%%%%%%%%%%%%%%
\section*{Main Convergence Theorem}
\addcontentsline{toc}{section}{Main Convergence Theorem}
\begin{theorem}[Ergodic $O(1/k)$ rate]\label{thm:ergodic}
Under Assumptions \ref{as:convex_smooth}–\ref{as:optimum} and with a stepsize $\eta\in(0,1/L]$, the iterates generated by \eqref{1} satisfy for all $K\ge1$
\[
F\bigl(\bar x_{K}\bigr)-F^\star
\le \frac{\|x_0-x^\star\|^{2}}{2\eta K},
\qquad 
\bar x_{K}:=\frac{1}{K}\sum_{k=0}^{K-1}x_{k},
\tag{4}
\]
where $x^\star\in\mathcal X^\star$ is any optimal point.
Consequently $F(\bar x_K)\to F^\star$ at the rate $O(1/K)$.
\end{theorem}

A non‑ergodic bound can also be derived:
\[
\min_{0\le k\le K-1}\bigl[F(x_k)-F^\star\bigr]\le \frac{\|x_0-x^\star\|^{2}}{2\eta K}.
\tag{5}
\]

%%%%%%%%%%%%%%%%%%%%%%%%%%%%%%%%%%%%%%%%%%%%%%%%%%%%%%%%%%%%%%%%%%%%%%%%%%%%
\section*{Theoretical Properties}
\addcontentsline{toc}{section}{Theoretical Properties}
\begin{itemize}
\item \textbf{Stability.} Because of the non‑expansiveness \eqref{3}, the mapping $x\mapsto\prox_{\eta g}(x-\eta\grad f(x))$ is \emph{firmly non‑expansive} when $\eta\le 1/L$, guaranteeing that the sequence $\{x_k\}$ remains bounded.
\item \textbf{Hyperparameter Sensitivity.} The stepsize must respect $\eta\le 1/L$. If a larger $\eta$ is used, the descent inequality \eqref{13} may fail and the method can diverge. Adaptive strategies (backtracking) can be employed to enforce an effective $\eta$ satisfying \eqref{2}.
\item \textbf{Comparison to Gradient Descent.} When $g\equiv0$, \eqref{1} reduces to standard gradient descent and \eqref{4} recovers the classic $O(1/k)$ rate for convex smooth functions. For non‑smooth $g$, the same rate holds without smoothing $g$.
\item \textbf{Extension to Strong Convexity.} If $F$ is $\mu$‑strongly convex, a linear rate $O\bigl((1-\eta\mu)^{k}\bigr)$ can be shown (see standard literature).
\end{itemize}

%%%%%%%%%%%%%%%%%%%%%%%%%%%%%%%%%%%%%%%%%%%%%%%%%%%%%%%%%%%%%%%%%%%%%%%%%%%%
\section*{Preliminaries (Definitions and Lemmas)}
\addcontentsline{toc}{section}{Preliminaries (Definitions and Lemmas)}
\begin{lemma}[Descent Lemma for $L$‑smooth $f$]\label{lem:descent}
For any $x,y\in\R^{d}$,
\[
f(y)\le f(x)+\langle\grad f(x),y-x\rangle+\frac{L}{2}\|y-x\|^{2}.
\tag{6}
\]
\end{lemma}
\begin{proof}
Standard proof using integral form of the gradient and Lipschitz continuity. $\square$
\end{proof}

\begin{lemma}[Proximal optimality condition]\label{lem:prox_opt}
For $z=\prox_{\eta g}(u)$,
\[
0\in \partial g(z)+\frac{1}{\eta}(z-u)
\quad\Longleftrightarrow\quad
\langle u-z,\,x-z\rangle\ge \eta\bigl[g(x)-g(z)\bigr],\;\forall x.
\tag{7}
\]
\end{lemma}
\begin{proof}
The definition of the proximal operator yields the inclusion. Multiplying by $\eta$ gives \eqref{7}. $\square$
\end{proof}

%%%%%%%%%%%%%%%%%%%%%%%%%%%%%%%%%%%%%%%%%%%%%%%%%%%%%%%%%%%%%%%%%%%%%%%%%%%%
\section*{Main Proof (Step‑by‑Step Derivation)}
\addcontentsline{toc}{section}{Main Proof (Step-by-Step Derivation)}
We prove Theorem \ref{thm:ergodic}.  Throughout we fix an arbitrary optimal point $x^\star\in\mathcal X^\star$.

\subsection*{Step 1: Apply the proximal optimality condition}
From \eqref{1} define $u_k:=x_k-\eta\grad f(x_k)$ and $x_{k+1}:=\prox_{\eta g}(u_k)$.  
By Lemma \ref{lem:prox_opt} with $z=x_{k+1}$, $u=u_k$, we have for any $x$,
\[
\langle u_k-x_{k+1},\,x-x_{k+1}\rangle\ge \eta\bigl[g(x)-g(x_{k+1})\bigr].
\tag{8}
\]

\subsection*{Step 2: Choose $x=x^\star$ in \eqref{8}}
\[
\langle u_k-x_{k+1},\,x^\star-x_{k+1}\rangle\ge \eta\bigl[g(x^\star)-g(x_{k+1})\bigr].
\tag{9}
\]

\subsection*{Step 3: Expand the inner product}
Recall $u_k=x_k-\eta\grad f(x_k)$.  Substituting into \eqref{9} gives
\[
\langle x_k-x_{k+1},\,x^\star-x_{k+1}\rangle
-\eta\langle\grad f(x_k),\,x^\star-x_{k+1}\rangle
\ge \eta\bigl[g(x^\star)-g(x_{k+1})\bigr].
\tag{10}
\]

\subsection*{Step 4: Relate the first term to a norm difference}
Using the identity $2\langle a,b\rangle = \|a\|^{2}+\|b\|^{2}-\|a-b\|^{2}$ with
$a=x_k-x_{k+1}$ and $b=x^\star-x_{k+1}$, we obtain
\[
2\langle x_k-x_{k+1},\,x^\star-x_{k+1}\rangle
= \|x_k-x^\star\|^{2}-\|x_{k+1}-x^\star\|^{2}-\|x_{k+1}-x_k\|^{2}.
\tag{11}
\]

Dividing \eqref{11} by $2$ and inserting into \eqref{10} yields
\[
\frac{1}{2}\Bigl(\|x_k-x^\star\|^{2}-\|x_{k+1}-x^\star\|^{2}-\|x_{k+1}-x_k\|^{2}\Bigr)
-\eta\langle\grad f(x_k),\,x^\star-x_{k+1}\rangle
\ge \eta\bigl[g(x^\star)-g(x_{k+1})\bigr].
\tag{12}
\]

\subsection*{Step 5: Rearrange terms}
Bring the gradient inner product to the right‑hand side:
\[
\frac{1}{2}\bigl(\|x_k-x^\star\|^{2}-\|x_{k+1}-x^\star\|^{2}\bigr)
-\frac{1}{2}\|x_{k+1}-x_k\|^{2}
\ge \eta\Bigl[\langle\grad f(x_k),\,x^\star-x_{k+1}\rangle
+g(x^\star)-g(x_{k+1})\Bigr].
\tag{13}
\]

\subsection*{Step 6: Insert the convexity of $f$}
Convexity of $f$ gives
\[
f(x^\star)\ge f(x_k)+\langle\grad f(x_k),\,x^\star-x_k\rangle.
\tag{14}
\]
Hence
\[
\langle\grad f(x_k),\,x^\star-x_{k+1}\rangle
= \langle\grad f(x_k),\,x^\star-x_k\rangle
+ \langle\grad f(x_k),\,x_k-x_{k+1}\rangle
\ge f(x^\star)-f(x_k)+\langle\grad f(x_k),\,x_k-x_{k+1}\rangle.
\tag{15}
\]

\subsection*{Step 7: Upper‑bound the last inner product using smoothness}
Apply Lemma \ref{lem:descent} with $x=x_k$, $y=x_{k+1}$:
\[
f(x_{k+1})\le f(x_k)+\langle\grad f(x_k),\,x_{k+1}-x_k\rangle
+\frac{L}{2}\|x_{k+1}-x_k\|^{2}.
\tag{16}
\]
Rearranging,
\[
\langle\grad f(x_k),\,x_k-x_{k+1}\rangle
\ge f(x_k)-f(x_{k+1})-\frac{L}{2}\|x_{k+1}-x_k\|^{2}.
\tag{17}
\]

\subsection*{Step 8: Substitute \eqref{15} and \eqref{17} into \eqref{13}}
\[
\begin{aligned}
\frac{1}{2}\bigl(\|x_k-x^\star\|^{2}-\|x_{k+1}-x^\star\|^{2}\bigr)
-\frac{1}{2}\|x_{k+1}-x_k\|^{2}
\;\ge\;&\eta\Bigl[f(x^\star)-f(x_k)\\
&\;+\;f(x_k)-f(x_{k+1})-\frac{L}{2}\|x_{k+1}-x_k\|^{2}\\
&\;+\;g(x^\star)-g(x_{k+1})\Bigr].
\end{aligned}
\tag{18}
\]

Simplify the right‑hand side:
\[
\eta\Bigl[F(x^\star)-F(x_{k+1})-\frac{L}{2}\|x_{k+1}-x_k\|^{2}\Bigr].
\tag{19}
\]

Thus we have
\[
\frac{1}{2}\bigl(\|x_k-x^\star\|^{2}-\|x_{k+1}-x^\star\|^{2}\bigr)
\ge \eta\bigl[F(x^\star)-F(x_{k+1})\bigr]
+\Bigl(\frac{1}{2}-\frac{\eta L}{2}\Bigr)\|x_{k+1}-x_k\|^{2}.
\tag{20}
\]

\subsection*{Step 9: Use the stepsize condition}
From Assumption \ref{as:stepsize}, $\eta L\le1$, therefore the coefficient
$\frac{1}{2}-\frac{\eta L}{2}\ge0$.  Dropping the non‑negative term yields the **fundamental descent inequality**
\[
\boxed{\;
\eta\bigl[F(x_{k+1})-F^\star\bigr]
\le \frac{1}{2}\bigl(\|x_k-x^\star\|^{2}-\|x_{k+1}-x^\star\|^{2}\bigr)
\;}
\tag{21}
\]
where $F^\star=F(x^\star)$.

\subsection*{Step 10: Telescoping sum}
Sum \eqref{21} for $k=0,\dots,K-1$:
\[
\eta\sum_{k=0}^{K-1}\bigl[F(x_{k+1})-F^\star\bigr]
\le \frac{1}{2}\bigl(\|x_0-x^\star\|^{2}-\|x_{K}-x^\star\|^{2}\bigr)
\le \frac{1}{2}\|x_0-x^\star\|^{2}.
\tag{22}
\]

Divide by $K$ and use convexity of $F$ (Jensen’s inequality) to obtain
\[
F\bigl(\bar x_{K}\bigr)-F^\star
\le \frac{1}{K}\sum_{k=1}^{K}\bigl[F(x_{k})-F^\star\bigr]
\le \frac{\|x_0-x^\star\|^{2}}{2\eta K}.
\tag{23}
\]

Equation \eqref{23} is exactly the bound \eqref{4} claimed in Theorem \ref{thm:ergodic}.
\qed

%%%%%%%%%%%%%%%%%%%%%%%%%%%%%%%%%%%%%%%%%%%%%%%%%%%%%%%%%%%%%%%%%%%%%%%%%%%%
\section*{Rate Derivation (With Constants)}
\addcontentsline{toc}{section}{Rate Derivation (With Constants)}
From \eqref{21} we also obtain a **non‑ergodic** bound.  For any $K\ge1$,
\[
F(x_{K})-F^\star
\le \frac{\|x_{K-1}-x^\star\|^{2}}{2\eta}
\le \frac{\|x_{0}-x^\star\|^{2}}{2\eta K},
\tag{24}
\]
where the second inequality follows by applying \eqref{21} recursively and discarding the non‑negative term $\|x_{k+1}-x_k\|^{2}$.  Hence
\[
\min_{0\le k\le K-1}\bigl[F(x_k)-F^\star\bigr]
\le \frac{\|x_{0}-x^\star\|^{2}}{2\eta K}.
\tag{25}
\]

Both \eqref{23} and \eqref{25} exhibit the $O(1/K)$ convergence rate with explicit constant $C=\|x_0-x^\star\|^{2}/(2\eta)$.

%%%%%%%%%%%%%%%%%%%%%%%%%%%%%%%%%%%%%%%%%%%%%%%%%%%%%%%%%%%%%%%%%%%%%%%%%%%%
\section*{Special Cases}
\addcontentsline{toc}{section}{Special Cases}
\begin{itemize}
\item \textbf{Pure Gradient Descent ($g\equiv0$).}  The proximal operator becomes the identity, and \eqref{21} reduces to the classic bound for gradient descent on smooth convex $f$.
\item \textbf{Indicator Function $g=\iota_{\mathcal C}$}.  Here $\prox_{\eta g}$ is the Euclidean projection onto a closed convex set $\mathcal C$.  The algorithm coincides with the projected gradient method, and the same $O(1/k)$ rate holds.
\item \textbf{Strongly Convex Composite $F$.}  If $F$ is $\mu$‑strongly convex, adding the inequality $\frac{\mu}{2}\|x_{k+1}-x^\star\|^{2}\le F(x_{k+1})-F^\star$ to \eqref{21} yields linear convergence:
\[
\|x_{k+1}-x^\star\|^{2}\le\bigl(1-2\eta\mu\bigr)\|x_k-x^\star\|^{2},
\]
provided $\eta\le 1/L$.
\end{itemize}

%%%%%%%%%%%%%%%%%%%%%%%%%%%%%%%%%%%%%%%%%%%%%%%%%%%%%%%%%%%%%%%%%%%%%%%%%%%%
\section*{Discussion}
\addcontentsline{toc}{section}{Discussion}
The proof above relies exclusively on three elementary ingredients:

\begin{enumerate}
\item \textbf{Lipschitz smoothness} of $f$ (Lemma \ref{lem:descent}),
\item \textbf{Convexity} of $g$ (proximal optimality condition Lemma \ref{lem:prox_opt}),
\item \textbf{Non‑expansiveness} of the proximal map (inequality \eqref{3}) which guarantees that the extra term $\|x_{k+1}-x_k\|^{2}$ appears with a non‑negative coefficient in \eqref{20}.
\end{enumerate}

Because each step was derived explicitly and every inequality was justified, the result holds under the minimal convexity assumptions stated.  The same template can be extended to stochastic proximal gradients, accelerated variants, or variable stepsizes by appropriately adapting Lemma \ref{lem:descent} and the proximal optimality condition.

\bigskip
\noindent\textbf{References for further reading}
\begin{itemize}
\item N. Parikh and S. Boyd, “Proximal Algorithms,” \emph{Foundations and Trends in Optimization}, 2014.
\item A. Beck and M. Teboulle, “A Fast Iterative Shrinkage‑Thresholding Algorithm,” \emph{SIAM J. Optim.}, 2009.
\end{itemize}

\end{document}